% Original diagnostic graph. Branches are hypotheses, not automatic diagnoses.
% Final width 112 mm; interventions return to knowledge/training decisions.
\begin{tikzpicture}[x=\linewidth,y=1mm,
 font=\figurefont\fontsize{8.5}{11}\selectfont,
 box/.style={draw=BookRule,fill=BookPaper,line width=.8pt,
   minimum height=11mm,align=center,inner sep=1.5mm},
 flow/.style={-{Latex[length=2mm]},draw=BookInk,line width=.9pt}]
 \node[box,text width=40mm] (loss) at (.5,0) {匹配预算后出现性能退化};
 \node[box,text width=61mm] (evidence) at (.5,-24)
   {补充诊断证据\\覆盖、容量、梯度冲突、模块替换};
 \node[box,text width=33mm] (training) at (.20,-53)
   {训练机制失效假设\\采样、表示或优化};
 \node[box,text width=33mm,fill=BookTeal!8] (structure) at (.80,-53)
   {共享结构失效假设\\动力学或接口不兼容};
 \node[box,text width=33mm] (trainfix) at (.20,-80)
   {修订训练方式\\重新做预算匹配比较};
 \node[box,text width=33mm] (structfix) at (.80,-80)
   {重估知识对象与范围\\重新设计任务划分};
 \draw[flow] (loss) -- (evidence);
 \draw[flow] (evidence.south) -- (.5,-38) -| (training.north);
 \draw[flow] (.5,-38) -| (structure.north);
 \draw[flow] (training) -- (trainfix);
 \draw[flow] (structure) -- (structfix);
 \node[anchor=north,text=BookMuted] at (.5,-91)
   {单次失败不足以确定根因；改动后重新验证。};
\end{tikzpicture}
